% ======================================================================
%  ResearchY-NP_170 — Lock-Lattice Phase Memory (Physical Realizability Test)
%  A proposal. Publication version of Docs/ResearchY/NP_NewPhysics/ResearchY-NP_170.md
%
%  Status: PROPOSAL — no experiment performed, no result claimed.
%  Scope:  no canonical AT claim is changed; no new primitive is introduced.
%          The locking nonlinearity and the engineered coupling network are
%          declared BOUNDARY imports (NP_005, NP_007/NP_011).
%
%  Build: pdflatex ResearchY-NP_170-Lock-Lattice-Program.tex   (twice)
% ======================================================================

\documentclass[11pt,a4paper]{article}

\usepackage[T1]{fontenc}
\usepackage[utf8]{inputenc}
\usepackage{lmodern}
\usepackage{amsmath,amssymb}
\usepackage{booktabs}
\usepackage{array}
\usepackage{enumitem}
\usepackage{xcolor}
\usepackage[hidelinks]{hyperref}
\usepackage[a4paper,margin=2.4cm]{geometry}

\newcommand{\ls}[1]{\ensuremath{\lambda_{#1}}}
\newcommand{\w}[1]{\ensuremath{\omega_{#1}}}
\newcommand{\Cninety}{\ensuremath{C_{96}(1..6)}}
\renewcommand{\arraystretch}{1.18}

\title{\textbf{Lock-Lattice Phase Memory}\\[0.35em]
\large A Physical Realizability Test for the Lock Law of\\[0.15em]
\large The Actualization Theory (AT)\\[0.35em]
\normalsize\emph{Proposal --- ResearchY-NP\_170}}

\author{Fabrice Wieser\\[1mm]
  \small Independent Researcher\\[0.5mm]
  \small Senior Software Engineer\\[0.5mm]
  \small \url{https://github.com/MagusDraconis/AT}}
\date{11 September 2026}

\begin{document}
\maketitle

\begin{abstract}
\noindent
The Actualization Theory (AT) --- a reconstruction of physics from Difference, Actualization and
Spectrum, in which every observable is derived from the D96 spectral structure while the Standard
Model Lagrangian is hosted rather than derived --- supports its lock law (Ch8; QG307--QG319), at
present, entirely by statistical evidence: datasets and synthetic cohorts. This document specifies an
experiment that would test it as a \emph{dynamical} claim, using an engineered nonlinear
oscillator ring whose coupling graph is the canonical circulant \Cninety. The proposal deliberately
abandons the energy-storage framing, which is refuted on density grounds, and identifies one idea
that survives: the lock lattice as a \emph{phase-memory} and \emph{threshold} element whose
distinguishing feature is a \emph{predicted} criticality rather than a fitted one. Three reasons
justify the experiment (theory falsification; phase memory where density is irrelevant; a
neuromorphic threshold element with predicted sharpness), and the exact procedure --- platform,
protocol, metrics, blinding, pre-registration and kill criteria --- is given in full. The
nonlinearity and the coupling network are declared as imported boundary inputs; the retention
hypothesis is flagged as an analogy across two AT systems, not a derivation. Both outcome branches
are informative, and the deciding experiment costs one bench prototype.
\end{abstract}

\section{Scope}

\paragraph{What AT is.} The Actualization Theory (AT) reconstructs physics on the canonical chain
\emph{Difference} $\to$ \emph{Actualization} $\to$ \emph{Inevitable Spectrum} $\to$
\emph{Physics}, from two primitives --- Difference (the fundamental boundary) and the tensor
reference $\eta$ --- with the observable attractor realized on the 96-site circulant
$C_{96}(1..6)$. The canonical monograph (v2.0, Ch1--Ch14) classifies every statement it makes as
DERIVED, EMERGENT or BOUNDARY. Throughout this proposal the acronym ``AT'' denotes that framework
only, and the phase labels AT-1xx / QGxxx denote its internal derivation and audit record.

This is a \textbf{proposal}. It contains no experimental result, changes no canonical AT claim,
reclassifies no value, and introduces no primitive. Two ingredients are imported and are declared as
boundary inputs: the locking nonlinearity (NP\_005: the locking force is ABSENT, a boundary) and the
engineered coupling network (NP\_007/NP\_011: the static coupling network is REFUTED as a physical
field). The experiment can fail to falsify AT; it cannot confirm it.

\section{The one idea worth doing}

\subsection{What is abandoned}
Energy storage. A 96-node lattice stores microjoule-scale energy at kilowatt-scale overhead if
cryogenic --- some $10^{6}$--$10^{9}$ times behind lithium-ion cells
($\sim 0.9~\mathrm{MJ\,kg^{-1}}$, $\sim 5000$ cycles, $\sim$100 EUR/kWh) on any storage figure of
merit. No amount of AT machinery changes that, so the storage claim is closed here.

\subsection{What survives}
The lock lattice as a \textbf{phase-memory element}: a device whose state variable is
\emph{which modes are locked}, rather than how much charge is separated (capacitor) or how much
chemical potential is stored (battery). Its retention hypothesis comes from the AT-120 one-way
barrier --- a topological charge, once actualized, cannot un-actualize --- and its read-out from the
AT-126 phase-interference law. We call it \textbf{LLPM} (Lock-Lattice Phase Memory):

\begin{quote}
\emph{An engineered circulant-$C_{96}(1..6)$ oscillator ring is charged into a locked mode
configuration; the locked configuration is predicted to retain its state longer than a matched
unlocked control, with the retention difference produced by the actualization barrier and readable
through $\Delta\theta$ interference.}
\end{quote}

The design is informative whether the answer is positive or negative, because AT's own
Ch14 concedes that ``standard complexity measures match or beat the lock rule on the evolving
cohort'': a statistical audit cannot settle the question. Only a physical phase transition can.

\section{Three reasons}

None of the three uses energy density as its metric.

\subsection{Reason 1 --- theory falsification (strongest)}
Every existing lock-law result (QG307--QG319) is computed on datasets or synthetic cohorts. A
hardware ring is the only place where Ch8 must produce a \emph{dynamical phase transition} in a
physical system instead of a statistic in a table.

\begin{table}[h]
\centering
\small
\begin{tabular}{@{}p{0.26\linewidth}p{0.66\linewidth}@{}}
\toprule
Item & Setting \\
\midrule
Metric & existence of a threshold in the \emph{physical} coupling, with sharpness $\ge 3$ and
width $\le 0.4$ (the QG316 criticality criteria) \\
Bar & hysteresis area $A>0$ at $\ge 5\sigma$ over $\ge 30$ runs, absent in all three controls \\
If it passes & the lock structure is physically realizable; the lock law stops being a candidate
descriptive statistic \\
If it fails & the lock law is demoted to a descriptive statistic (consistent with QG312's partial
fake and QG319's $96.9\%$ false-positive / $66.1\%$ false-negative rates), and the storage claim is
withdrawn permanently \\
\bottomrule
\end{tabular}
\end{table}

\noindent Both branches are publication-grade information, at the cost of one bench prototype.

\subsection{Reason 2 --- phase / coherence memory}
Here density is irrelevant; the metric is energy per operation: write energy, retention, read-out
fidelity, variability. Because the state variable is \emph{phase}, not charge, writing is
phase-pinning (M\_002) rather than charge transfer, so the write-energy floor is set by phase noise
rather than by Coulomb energy. The bar to matter is $<{\sim}10$~fJ per write, retention comparable
to a supercapacitor's self-discharge window, and variability $<5\%$ (competitive with
MRAM/PCM/FeRAM). The failure mode is thermal and technical phase diffusion: a room-temperature ring
may retain nothing beyond $1/\omega_0 \times Q$.

\subsection{Reason 3 --- neuromorphic threshold element}
The lock transition is a sigmoidal threshold with hysteresis --- a synapse/neuron primitive. What is
AT-specific is that the threshold position and sharpness are \emph{predicted} by the operator and
lock structure rather than fitted, which is precisely what phase-change and memristive devices
cannot offer. The bar: beat PCM/memristor variability and switching energy at equal node count. The
risk: the nonlinearity that produces locking is imported (NP\_005) and may itself dominate the
variability.

\section{Canonical basis: the numbers the hardware must realize}

The observable attractor is the circulant \Cninety{} (degree 12) with Laplacian spectrum

\begin{equation}
\ls{k} \;=\; 2\sum_{d=1}^{6}\Bigl(1-\cos\frac{2\pi dk}{96}\Bigr),
\qquad k=0,\dots,95 .
\label{eq:spec}
\end{equation}

Equation~\eqref{eq:spec} was recomputed independently for this proposal; the checks are listed in
Table~\ref{tab:numbers}.

\begin{table}[h]
\centering
\small
\begin{tabular}{@{}lccc@{}}
\toprule
Quantity & Recomputed & Canonical & Check \\
\midrule
$\ls{2}$ (first positive eigenvalue, $k=1$) & $0.386351$ & $0.3864$ & \checkmark \\
$\w{1}=\sqrt{\ls{2}}$ & $0.621571$ & $0.6216$ & \checkmark \\
$\w{\max}$ ($k=48$, $\lambda=12$) & $3.979620$ & --- & --- \\
$\mathrm{span}=\w{\max}/\w{1}$ & $6.402515$ & $6.4025$ & \checkmark \\
octave bands over the 95 positive modes & $[4,4,87]$ & $[4,4,87]$ & \checkmark \\
zero modes & $1$ & $1$ & \checkmark \\
mirror pairs $\ls{k}=\ls{96-k}$ & 47 pairs $+$ self-conjugate $k=48$ & 47 pairs & \checkmark \\
multiplicity multiset & $1+42\times 2+5+6=95$ & $[42\times2,5,6]$ & \checkmark \\
distinct eigenvalues & $45$ & --- & --- \\
lock chain (D96-specific) & $\mathrm{occMom}/\Sigma m = 20.0026$ & $20.0026$ & Ch8 \\
\bottomrule
\end{tabular}
\caption{Verified spectrum data for \Cninety. The canonical lock values are listed for reference
only: by QG313 the lock \emph{structure} is universal while the lock \emph{values} are
domain-specific, so a hardware ring must derive its own lock values and must not be scored against
$20.0026$.}
\label{tab:numbers}
\end{table}

Two canonical facts govern what hardware may claim. \textbf{(i) QG313:} lock structure is universal,
lock values are domain-specific; the D96 values ($20.0026$, $2.4105$, $8.2980$) are therefore not
hardware predictions. \textbf{(ii) QG316:} the critical transition $g^\ast\approx 0.31$ is a
\emph{synthetic-cohort} parameter --- an exponent fraction over a 40-step organization ramp,
disclosed as synthetic by MONO007 A08 --- not a physical coupling, and it is not transferable. The
hardware prediction is the \emph{existence} of a threshold satisfying the criticality criteria, not
the number $0.31$. The inverse-spectral construction that converts a target spectrum into coupling
weights already exists as an accepted ResearchY result (T\_001, inverse DFT onto the first
circulant row); no new mathematics is required to specify the hardware.

\section{Hypotheses and pass/fail criteria}

\begin{table}[h]
\centering
\small
\begin{tabular}{@{}p{0.05\linewidth}p{0.24\linewidth}p{0.22\linewidth}p{0.41\linewidth}@{}}
\toprule
\# & Hypothesis & Pass & Fail / interpretation \\
\midrule
H0 & the ring realizes \Cninety & mirror-pair error $<1\%$; bands $[4,4,87]$
& absence $=$ \emph{validity} failure, not a theory failure \\
H1 & a lock threshold exists in the physical coupling & sharpness $\ge 3$ and width $\le 0.4$ in
$f(g)$, hysteresis $A>0$ at $\ge5\sigma$ & $f(g)$ smooth and reversible $\Rightarrow$ lock law
demoted to descriptive statistic \\
H2 & locked configurations retain state longer than controls & two time scales,
$R=\tau_{\rm slow}/\tau_{\rm fast}\ge 10$, single-exponential in all controls & same decay law as
controls $\Rightarrow$ the AT-120 analogy does not transfer; retention claim withdrawn \\
H3 & read-out follows the interference law & $|\Theta_{\rm mid}|\propto\cos(\Delta\phi/2)$,
visibility $V\ge0.90$ & $V<0.90$ or wrong functional form \\
H4 & threshold element is competitive & beats PCM/memristor switching energy and variability &
reason 3 retired; reasons 1--2 unaffected \\
\bottomrule
\end{tabular}
\caption{Hypotheses. H2 is flagged as the highest-risk, highest-information hypothesis: it is an
analogy across two different AT systems (the $R$-field condensate of AT-120 and the circulant
Laplacian ring of Ch6--Ch8), not a derivation. Testing an analogy is legitimate; asserting it as
derived would not be.}
\label{tab:hyp}
\end{table}

\section{Apparatus}

\subsection{Tier 1 (primary; bench electronic, days, low cost)}
A 96-node ring of nonlinear oscillators with programmable circulant coupling.

\begin{table}[h]
\centering
\small
\begin{tabular}{@{}p{0.24\linewidth}p{0.68\linewidth}@{}}
\toprule
Element & Specification \\
\midrule
Nodes & 96 active resonators (LC tank $+$ saturating amplifier), $f_0=100$~kHz--1~MHz \\
Nonlinearity & $\tanh$-type saturation (Adler/Kuramoto phase dynamics) --- \emph{imported, BOUNDARY
(NP\_005)} \\
Coupling & for node $i$ and each offset $d\in\{1..6\}$: links to $i\pm d$ via voltage-controlled
transconductance $g_{m,d}$; 12 links per node \\
Programmability & one DAC per offset sets $w_d$; the Laplacian is then circulant with
$\ls{k}=2\sum_d w_d(1-\cos 2\pi dk/96)$ --- the same family as Eq.~\eqref{eq:spec}, with weights
from T\_001 \\
Sweep & $g\equiv\sum_d w_d / \sum_d w_{d,\max}$, scanned over $g\in[0,1]$ in 40 steps (mirroring the
QG316 ramp) \\
Controls & (a) detuned ring (same $Q$, $\w{k}$ shifted out of ratio); (b) random-coupling ring (same
$Q$ and span); (c) linear-ramp spectrum (must fail CROWDING per QG305/QG308) \\
Read-out & 96-channel multiplexed IQ demodulation at $f_0$, $\ge1$~MS/s per channel; $\theta_j$ and
$A_j$ per node \\
Injector & single-channel drive at selectable $\w{k}$ with programmable phase $\theta_{\rm in}$;
abruptly removed for the retention run \\
Budget & 2--4 $\times$ 16-channel ADC/DAC plus FPGA for real-time IQ: the EUR~5--20k class \\
\bottomrule
\end{tabular}
\end{table}

\noindent Design constraints: node-to-node $f_0$ spread must be $<1\%$ or the H0 test fails for
trivial reasons; phase noise sets the retention floor, since the locked/unlocked contrast is only
measurable while $1/(Q\w{0})\ll\tau_{\rm slow}$; the nonlinearity strength is the analogue of the
QG316 organization parameter, so the mapping must be pre-registered before data-taking and the raw
control voltage reported alongside the normalized $g$.

\subsection{Tier 2 (photonic chip, months)}
96 coupled SiN microring resonators; the transmitted spectrum gives $\w{k}$ directly (H0), locking
via the Kerr nonlinearity, read-out by phase-sensitive heterodyne detection. Costlier, better for
H1, worse for H2 (loss).

\subsection{Tier 3 (superconducting, ${\sim}1$ year)}
96 coupled superconducting cavities/transmons at $T<20$~mK, $Q\approx10^{6}$: the only tier in which
a long retention claim is meaningful. Justified only if Tier~1 yields $A>0$ at $\ge5\sigma$,
$R\ge10$, and a device-level figure of merit beating a supercapacitor. Cryogenics are never
justified by the energy stored.

\section{Protocol}

All runs use $\ge30$ independent repetitions per configuration, with the analysis script frozen
before the first run.

\begin{description}[leftmargin=1.2em,style=nextline]
\item[R0 --- calibration] measure each node's free-running $\w{0}$ and $Q$; trim or calibrate to
$<1\%$ spread; record the spread as a systematic.
\item[R1 --- validity (H0)] set $w_d$ to the \Cninety{} values; measure $\w{k}$ by ring-down or
spectral scan; compute the mirror-pair error and the octave band occupancy (expect $[4,4,87]$ over
95 positive modes). \emph{Gate:} if this fails, stop --- the apparatus is not a D96 lattice and
nothing downstream is interpretable.
\item[R2 --- charging and read-out (H3)] drive at $\w{k}$ to settling, then probe two adjacent sites
with controlled relative phase $\Delta\phi$; measure
$|\Theta_{\rm mid}|/(|\Theta_{x_1}|+|\Theta_{x_2}|)$ versus $\Delta\phi$ and fit $\cos(\Delta\phi/2)$.
\emph{Gate:} $V\ge0.90$.
\item[R3 --- lock threshold and hysteresis (H1)] sweep $g$ upward over 40 steps and then downward;
at each step, after settling, compute
\[
f(g)=\frac{\#\{j:\;|d\Delta\theta_j/dt|<\varepsilon\}}{96},\quad \varepsilon=0.01\,\w{1},
\qquad
A=\int|f_{\rm up}(g)-f_{\rm down}(g)|\,dg ,
\]
with $\text{sharpness}=\max|\Delta f|/\overline{|\Delta f|}$ and width $=$ the $g_{10}\!\to\!g_{90}$
interval. Test sharpness $\ge3$, width $\le0.4$, $A>0$ at $\ge5\sigma$. Report the measured
$g_c$ as a hardware-specific value (QG313) and do \emph{not} compare it with $0.31$.
\item[R4 --- retention (H2)] (i) \emph{ring-down}: charge into a locked configuration, remove the
drive abruptly, record total stored energy $\sum_j A_j^2/2$ versus $t$ for $\ge10^4/\w{1}$;
(ii) \emph{cycle test}: $\ge10^3$ charge/hold/discharge cycles, recording state fidelity after hold.
Fit
\[
\text{locked: } E(t)=a\,e^{-t/\tau_{\rm fast}}+b\,e^{-t/\tau_{\rm slow}},
\qquad
\text{control: } E(t)=a\,e^{-t/\tau_c},
\]
and report $R=\tau_{\rm slow}/\tau_{\rm fast}$ and the storage figure of merit
$\mathrm{FOM}=R\cdot E_{\rm ret}/E_{\rm in}$. \emph{Pass:} $R\ge10$ with the two-time-scale fit
preferred by BIC/AIC against the single-exponential null, with all three controls
single-exponential.
\item[R5 --- threshold benchmark (H4, optional)] only if R3 passes: measure switching energy per
node, threshold sharpness and board-to-board variability across five identical boards; compare with
published PCM/memristor figures.
\item[R6 --- controls and blinding] labels $\{$locked, control a, b, c$\}$ are hidden from the
analyst until the frozen script has produced numbers (QG319: FP $96.9\%$, FN $66.1\%$ --- precisely
the failure mode to guard against). The lock classifier itself is fixed before R3.
\end{description}

\section{Statistics, blinding and reproducibility}

\begin{table}[h]
\centering
\small
\begin{tabular}{@{}p{0.28\linewidth}p{0.64\linewidth}@{}}
\toprule
Rule & Setting \\
\midrule
Repeats & $\ge30$ per configuration; report mean $\pm$ 95\% CI and effect size, never a single trace \\
Significance & $A>0$ and $R\ge10$ each at $\ge5\sigma$ \\
Model comparison & BIC and AIC for single- vs two-time-scale decay; the null is always the simpler
model \\
Blinding & label blinding with a scripted or third-party unblinding step \\
Freezing & analysis-script hash recorded before R3; any change restarts the run \\
Negative controls & detuned, random-coupling and linear-ramp lattices must fail H1 and H2 \\
Systematics & node-frequency spread, phase noise, ADC clock jitter, temperature drift and cable
delay budgeted explicitly \\
Reproducibility & raw IQ traces archived with instrument settings and seeds \\
\bottomrule
\end{tabular}
\end{table}

\section{Kill criteria and decision tree}

\begin{quote}\ttfamily\small
R1 fails \hfill apparatus problem $\to$ repair and rerun (not a theory result)\\
R1 passes, R3 null \hfill H1 falsified $\to$ withdraw the storage claim permanently\\
R3 passes, R4 null \hfill keep the threshold-device direction (R3), drop phase memory (R2)\\
R3+R4 pass, FOM $<10\times$ supercap \hfill stop; do not escalate on storage grounds\\
R3+R4 pass, FOM $\ge10\times$ supercap \hfill escalate to Tier 2, then Tier 3
\end{quote}

\noindent Cryogenics are never justified by the energy stored.

\section{Boundaries: what is imported}

\begin{table}[h]
\centering
\small
\begin{tabular}{@{}p{0.30\linewidth}p{0.30\linewidth}p{0.34\linewidth}@{}}
\toprule
Imported & Status & Consequence \\
\midrule
locking nonlinearity & BOUNDARY (NP\_005: locking force ABSENT) & the experiment tests the lock
\emph{structure}, not a derived interaction \\
engineered coupling network & BOUNDARY (NP\_007/NP\_011: REFUTED as a physical field) & the lattice
is built, never claimed emergent \\
retention mechanism (AT-120 barrier) & analogy across systems & H2 is a hypothesis about
transferability, not a derived prediction \\
$g^\ast=0.31$ & synthetic-cohort parameter (QG316; MONO007 A08) & a hardware $g_c$ is expected to
differ; quoting $0.31$ as a hardware number would be a fit \\
D96 lock values $20.0026/2.4105/8.2980$ & D96-specific (QG313) & derive the hardware's own lock
values from its own spectrum \\
energy-storage claims & out of scope & the density argument is accepted as fatal \\
\bottomrule
\end{tabular}
\end{table}

\noindent This proposal alters the DERIVED / EMERGENT / BOUNDARY status of no existing result.

\section{What each outcome means}

\begin{table}[h]
\centering
\small
\begin{tabular}{@{}p{0.24\linewidth}p{0.36\linewidth}p{0.34\linewidth}@{}}
\toprule
Outcome & Meaning for AT & Meaning for engineering \\
\midrule
R1 only & the D96 lattice is realizable & an oscillator-ring testbed with a derivable spectrum \\
R1+R3 & the lock structure is physically realizable, not merely statistical & a threshold device
with predicted criticality criteria \\
R1+R3+R4 & phase memory with barrier-protected retention & a candidate phase/coherence memory
element (reason 2) \\
R1+R4 without R3 & two-time-scale decay without a threshold --- suspects an ordinary high-$Q$ effect
& no device claim; the lock reading is unsupported \\
R1 only, R3 null & lock law demoted to a descriptive statistic & the platform stays useful; the
storage idea ends \\
\bottomrule
\end{tabular}
\end{table}

\noindent The negative branch is not a failure of the program: it is the honest closure of an open
question, and it is the branch AT's own Ch14 anticipates.

\section{Deliverables}

\begin{enumerate}[leftmargin=1.4em]
\item This proposal (PLANNED).
\item A numerical precursor (planned): a deterministic Kuramoto/Adler ring simulation on
\Cninety{} (fixed seeds, no randomness) computing the expected sharpness/width curves and the
two-time-scale signature, so that R3/R4 have an \emph{a priori} prediction.
\item A canonical prediction-registry entry, quoted below, corresponding to \textbf{AT-P043}
(AT-P042 being the tick discriminator, M\_009). The canonical registry is \emph{not} modified by
this proposal.
\end{enumerate}

\begin{quote}\raggedright\scriptsize\ttfamily
\string\predstart\{AT-P043\}\{PROPOSED\}\\
\string\pfield\{Title\}\{Lock-lattice physical realizability: a lock threshold in an engineered C96
ring.\}\\
\string\pfield\{Derived From\}\{Difference $\to$ Actualization $\to$ Spectrum (D96) $\to$ Locking
operator $\to$ lock law (Ch8); criticality criteria from QG316; design recipe from ResearchY-T\_001.\}\\
\string\pfield\{Prediction\}\{An engineered circulant C96(1..6) nonlinear oscillator ring exhibits a
lock threshold in the physical coupling (sharpness $\ge$ 3, width $\le$ 0.4) with hysteresis, and a
locked configuration decays with two time scales against a single-exponential control. Structure
universal (QG313); the threshold value and the lock values are hardware-specific and are NOT
predicted.\}\\
\string\pfield\{Experimental Test\}\{Bench electronic 96-node ring, runs R1--R4, $\ge$30 repeats,
blinded labels, three negative controls.\}\\
\string\pfield\{Current Evidence\}\{None --- proposal. All existing lock-law evidence is statistical
(QG307--QG319).\}\\
\string\pfield\{Falsification Criterion\}\{No hysteresis, or locked retention decay
indistinguishable from all three controls.\}
\end{quote}

\section{Conclusion}

One idea survives the density objection: the lock lattice as a phase-memory and threshold element
whose distinguishing feature is a predicted criticality, tested on hardware for the first time.
Three reasons justify it --- falsification of the lock law as a dynamical claim (strongest), phase
memory where density is irrelevant, and a neuromorphic threshold element with predicted sharpness.
The deciding experiment costs one bench prototype, and both branches of the outcome are informative.
Energy storage is explicitly abandoned, the nonlinearity and coupling network are explicitly
imported, and the retention hypothesis is explicitly an analogy awaiting test rather than a
derivation.

\bigskip
\noindent\textbf{Status:} PROPOSAL --- no experiment performed, no result claimed, no canonical AT
statement modified.

\end{document}
